\chapter{Návrh ovladačů}

Navržený obvod (viz přílohy C a D) představuje digitální vstupní rozhraní určené pro snímání stavů
osmi mechanických spínačů a jejich sériový přenos do mikrokontroléru prostřednictvím minimálního
počtu signálových vodičů.

Vstup z tlačítek je veden přes posuvný registr CD4021, který realizuje paralelní vstup a sériový
výstup (posuvný registr typu PISO -- \textit{Parallel-In Serial-Out}). Ten je řízen třemi signály:

\begin{itemize}
    \item \textbf{STROBE} -- aktivní hrana způsobí paralelní načtení okamžitých stavů všech vstupů do
        interního registru,
    \item \textbf{CLK} -- hodinový signál synchronizující postupné vysouvání dat na sériový výstup,
    \item \textbf{OUT} -- sériový datový výstup.
\end{itemize}

Protokol přenosu je striktně sekvenční: puls signálu STROBE uzamkne aktuální stav vstupů, načež
každá náběžná hrana CLK vysune nejstarší uložený bit na výstup OUT. Stav osmi tlačítek je tak
přenesen v~osmi hodinových cyklech, aniž by bylo nutné rozšiřovat sběrnici o~další datové vodiče.

Každý spínač je zapojen v konfiguraci \textit{active-low} -- jeho stisknutí je považováno za
logickou 0 (napětí blízké referenčními GND).

Potlačení zákmitů tlačítek je řešeno skutečností, že čtení vstupu probíhá striktně před začátkem
emulace snímku. Je pravděpodobný, že v okamžik čtení stavu tlačítek zákmit už není přítomen. V
nejhorším případě, když čtení proběhne v přítomnosti zákmitu, může to způsobit zpoždění vstupu, ale
v praxi je to zanedbatelný. Navíc je zákmit potlačen přirozenou délkou doby snímání dané hodinovou
frekvencí CLK.
